% ch3_literature.tex — 第三章 文献综述
\section{文献综述}

\subsection{在线评论研究的两类视角}

在线评论研究存在两条平行的研究路径——评论对潜在购买者的影响研究与评论者自身心理变量关系研究，两者在研究对象、理论视角和方法论上存在根本差异。

第一类研究聚焦评论作为信息源对潜在购买者决策的影响机制。Aridor et al.\cite{aridor2022}通过实地实验证实在线推荐通过塑造消费者信念影响选择，信念强度对决策权重的影响在低经验消费者中更显著。Su et al.\cite{su2022}分析中国电商平台大规模评论数据，验证情感倾向对决策的预测作用。Cao et al.\cite{cao2021}的元分析整合近十年eWOM研究，证实eWOM对购买意向的正向影响受平台类型和产品类别调节。Filieri et al.\cite{filieri2022}通过TripAdvisor实证研究发现评论可信度与有用性是影响消费者采纳决策的关键因素。在评论有用性研究方面，Jiménez \& Díaz\cite{jimenez2022}、Baka et al.\cite{baka2022}、Duan et al.\cite{duan2023}和Liu et al.\cite{liu2022}的研究均表明评论者特征、产品类别和平台差异对评论有用性产生系统性影响。

第二类研究聚焦评论者自身在评论中表达的心理变量及其关系。已验证购买的评论者在评论中同时表达基于实际使用体验形成的产品态度与未来行为意向（如回购、推荐），这种"态度--行为意向"关系是理解消费者忠诚度和口碑传播的重要视角。根据计划行为理论\cite{ajzen1991}，态度是行为意向的重要前因变量。

现有研究在理论视角和方法论上存在两个显著局限。第一，绝大多数研究聚焦第一类视角，而忽视评论者自身心理变量关系的检验，导致对评论生成机制的理解不足。第二，现有态度--行为意向关系研究主要依赖问卷调查数据，较少从自然语言文本中提取心理变量并检验其关系，限制了研究的生态效度和可扩展性。本研究聚焦第二类视角，利用大语言模型从评论文本中提取态度和行为意向变量，检验态度--行为意向关系，填补这一研究空白。

表\ref{tab:two_paradigms}对两类研究路径作出系统对比，明确本研究的定位。

\begin{table}[htbp]
\centering
\caption{在线评论研究两类范式对比}
\label{tab:two_paradigms}
\begin{threeparttable}
\begin{tabular}{@{}lll@{}}
\toprule
对比维度 & 第一类：评论$\to$他人决策 & 第二类：本研究定位 \\
\midrule
研究对象 & 潜在购买者 & 评论者自身 \\
核心变量关系 & 评论特征$\to$消费者采纳/购买 & 信念$\to$态度$\to$评分/行为意向 \\
数据来源 & 问卷、实验、平台聚合数据 & 评论文本（自然发生数据） \\
方法论工具 & 传统实证方法 & LLM文本提取＋平台客观数据 \\
CMV控制 & 依赖程序控制或实验设计 & 三独立数据源从设计层面消除 \\
\bottomrule
\end{tabular}
\begin{tablenotes}
\small
\item 注：CMV = Common Method Variance，同源方差。
\end{tablenotes}
\end{threeparttable}
\end{table}

\subsection{消费者信念与态度的实证研究}

消费者信念是态度形成的基础单元，信念的识别、评价与强度评估构成多属性态度测量的核心操作化问题。

Aridor et al.\cite{aridor2022}的实地实验证实在线推荐通过塑造消费者信念影响购买选择，信念强度对决策权重的影响在低经验消费者中更显著，为信念强度测量的重要性提供了实证依据。Su et al.\cite{su2022}分析大规模评论数据发现情感倾向对决策的预测作用，表明信念评价（情感极性）是态度测量的关键维度。在信念提取的粒度选择上，面向级情感分析（ABSA）理论强调对具体属性的情感提取以支持细粒度分析\cite{zhang2022absa}，但同时承认整体评价在理解用户态度中的价值：当评论缺乏明确属性表达时，提取整体评价性信念有助于更完整地捕捉消费者总体感受。在信念强度评估方面，Sun et al.\cite{sun2023}的研究表明基于LLM的方面级情感分析在处理情感强度评估时能综合考虑词汇语义、上下文语境和否定结构，准确率显著优于传统词典方法。Yin et al.\cite{yin2023}指出在线评论存在系统性积极偏向，但这一偏向受平台设计和消费者心理因素共同调节，语境对情感词汇解读至关重要。

现有研究存在三方面局限：信念提取研究多聚焦属性级情感分析的技术实现，较少讨论如何在Fishbein多属性态度理论框架下操作化信念概念；信念强度评估缺乏将其作为连续变量纳入态度计算公式的实证检验；现有研究多采用传统NLP方法，较少探索大语言模型在完整Fishbein框架信念提取中的应用潜力。本研究基于Fishbein多属性态度理论设计LLM提示词，从评论文本中提取信念陈述、情感极性（$e_i$）和信念强度（$b_i$），通过加权求和计算多属性态度（ATT），为态度测量提供新的操作化路径。

\subsection{行为意向测量}

行为意向（Behavioral Intention）是指消费者执行特定行为的倾向和计划，是态度影响消费者决策的重要结果变量\cite{fishbein1975}。Aridor et al.\cite{aridor2022}的研究表明，推荐系统通过影响消费者的信念来驱动消费行为，行为意向测量需要考虑信念强度的影响。

行为意向的测量通常采用量表形式，根据Lim et al.\cite{lim2022}的研究综述和Ajzen\cite{ajzen1991}的计划行为理论，行为意向是预测实际行为的最直接前因变量。Sheeran \& Webb\cite{sheeran2022}的元分析研究表明，行为意向与实际行为之间存在中等强度的相关关系，且这一关系受情境因素的调节。本研究采用四级量表，区分不同强度的行为意向，以更准确地捕捉消费者未来的行为倾向。

在电商情境中，行为意向通常包括再次购买意向（repurchase intention）和推荐意向（recommendation intention）。值得注意的是，在线评论文本中表达的行为意向属于"自愿披露的意向"，消费者选择在评论中表达未来意向是一种主动行为，因而可能并非随机样本——这构成了本研究H3/H4分析需要诚实讨论的选择偏差问题（详见第四章）。

\subsection{评论评分Rating作为态度外显指标}

评论评分（Rating）是消费者在完成购买并使用产品后，通过平台提交的整体满意度数值评价，通常以星级形式（1--5星）呈现。与评论文本不同，Rating来自平台结构化字段，不受消费者文本表达能力的限制，能够直接反映消费者对产品的整体态度判断，是消费者态度的客观外显表达（objective behavioral indicator of attitude）。

Duan et al.（2023）\cite{duan2023}的研究指出，在线评分数据普遍呈现显著的J型分布（J-shaped distribution），即评分极端值（5星和1星）占比远高于中间值，这一分布模式在Amazon等主流电商平台上尤为明显。这种评分分布的偏态特征在统计分析上具有重要含义：若直接使用OLS线性回归，将违反误差正态分布假设，导致系数估计有偏；Kim et al.（2022）\cite{kim2022}的研究也指出，忽视评分的非正态分布会产生误导性的统计推断，因此有序Logistic回归或Tobit模型是处理评分数据的更严谨选择。

Rating与评论文本的关系也是重要研究议题。Li et al.（2022）\cite{li2022}在对亚马逊等平台评论有用性的系统综述中指出，Rating作为评论信息诊断性（review diagnosticity）的组成部分，与评论有用性（helpful\_vote）共同构成消费者决策信息的来源——这与本研究将helpful\_vote作为控制变量的设计形成理论呼应。此外，Yin et al.（2023）\cite{yin2023}的研究表明，在线评分存在系统性正向偏差（positivity bias），消费者倾向于在公开平台上给出高于其私人评价的评分，这进一步说明Rating与消费者内在态度之间存在系统性偏差，而从评论文本中提取的多属性态度（ATT）可能能够捕捉这种Rating所掩盖的复杂性。

综上，本研究将Rating定位为"消费者态度的客观外显表达"，而非等同于消费者内在态度本身；ATT（从评论文本中LLM提取的多属性态度）与Rating（平台客观评分）形成"多方法测量同一构念"的格局，两者测量方法的完全独立性正是H1/H2分析避免CMV问题的根本依据。

\subsection{评论者经验与消费者专业性}

评论者经验（Reviewer Experience）是本研究引入的调节变量，旨在捕捉不同评论者在产品类别熟悉度和消费者专业性方面的差异。相关理论基础来自消费者行为研究中对"消费者专业性"（consumer expertise）的系统研究。

Mitchell et al.（2022）\cite{mitchell2022}在多领域实证分析中将消费者专业性定义为通过实践积累的主观能力和经验，包含认知努力、认知结构、分析能力、精细化能力和记忆五个维度。其中，"品类熟悉度"（category familiarity）是专业性最基础的维度，指消费者在特定品类中的经验积累程度。本研究以用户在Amazon Beauty品类的历史评论总数衡量评论者经验，代理其品类熟悉度。

Rocklage et al.（2023）\cite{rocklage2023}的研究从另一角度支持引入评论者经验作为调节变量：个体通过直接体验形成的态度，相较于通过间接信息（如广告、他人推荐）形成的态度，具有更高的可及性（attitude accessibility）和清晰度，因而与外显行为之间的一致性显著更强。在本研究情境中，评论者经验越丰富，其对产品属性的判断越基于直接体验积累，形成的多属性态度越清晰，与评分行为的一致性应越强。

需要说明的是，评论数量作为消费者专业性的代理指标存在一定局限性。Amazon数据集不提供用户性别、年龄、购买金额等人口统计信息，评论数量是可获取的最接近"品类熟悉度"的客观指标。此外，消费者专业性对态度--行为一致性的影响方向并非无争议：高经验评论者在积累丰富评价经验的同时，也可能形成跨产品的比较性评价框架，使其评分行为不再单纯映射当前产品态度，而是融入竞品参照、性价比权衡等外部因素；这种复杂化评价机制可能使高经验评论者的ATT与Rating之间的对应关系趋于松弛而非收紧。因此，评论者经验对ATT--Rating关系的调节方向是一个有待实证检验的开放性问题，本研究将H2设定为探索性假设，不对方向作出预判。这一测量局限将在研究局限性部分详细讨论。

\subsection{大语言模型在文本分析中的应用}

大语言模型（LLM）技术的快速发展为从非结构化文本中提取结构化心理变量提供了新的方法论工具，但其在消费者行为研究中的应用仍处于探索阶段。

Zhao et al.\cite{zhao2023}对主流大语言模型的系统综述表明，LLM在复杂推理、文本生成和信息提取等任务上展现出前所未有的能力，为社会科学研究提供了新的方法论工具。Malik \& Bilal（2024）\cite{malik2024}的系统综述表明，在线评论分析领域正从规则方法向深度学习和大语言模型范式快速迁移。Li et al.（2025）\cite{li2025llm}提出的"LLM-as-a-judge"范式为利用LLM对文本内容进行结构化评估提供了方法论基础，研究表明经过适当提示设计的LLM在评分任务上可达到与人类专家较高的一致性。Grimmer et al.（2022）\cite{grimmer2022}在《Text as Data》中系统阐述了将文本数据化方法应用于社会科学研究的框架，强调文本分析方法应服务于研究的理论问题而非追求技术复杂性。Egami et al.（2022）\cite{egami2022}进一步探讨了如何利用文本数据进行因果推断，为文本分析在社会科学研究中的应用提供了方法论指导。

现有研究存在三方面局限：LLM在消费者行为研究中的应用多聚焦情感分析等描述性任务，较少探索如何提取符合心理学理论定义的潜变量；LLM方法的可靠性验证研究不足，多数研究仅报告准确率等单一指标，缺乏系统的信度和效度检验；LLM提取结果与传统测量方法的关系尚不明确，限制了其在理论检验中的应用。本研究采用DeepSeek大语言模型，通过精心设计的提示词从评论文本中提取消费者信念、信念强度和行为意向，并系统验证LLM提取方法的预测效度（H1）和增量效度（H4），为LLM方法在消费者行为研究中的应用提供方法论范式。
